% uid: ""
% school: "lycee"
% level: "1ere"
% chapter: "Théorème de Thalès et sa réciproque"
% title: "La charpente"
% difficulty: "Moyen"
% tags: ["thales", "segments-paralleles", "somme-de-termes", "figure"]
% provenance: "own"
% batch: "2026-11-22-revision-1ere-A"

% ---STATEMENT---
Le toit d'un hangar est un triangle $ABC$ avec $BC=12$ m et une hauteur $AH=6$ m ($H\in[BC]$). On fixe cinq barres horizontales $[M_kN_k]$ ($k=1,2,3,4,5$), avec $M_k\in[AB]$, $N_k\in[AC]$ et $(M_kN_k)\parallel(BC)$, de sorte que les points $M_k$ partagent $[AB]$ en six parties égales : $AM_k=\dfrac k6\,AB$.
% alt: Triangle ABC de base 12 m avec cinq barres horizontales M_kN_k qui partagent AB en six parties égales
\begin{center}
\begin{tikzpicture}[scale=0.62]
  \coordinate (B) at (0,0); \coordinate (C) at (12,0); \coordinate (A) at (5,6); \coordinate (H) at (5,0);
  \draw[thick] (A)--(B)--(C)--cycle; \draw[dashed] (A)--(H);
  \foreach \k in {1,...,5}{
    \coordinate (M\k) at ($(A)!{\k/6}!(B)$); \coordinate (N\k) at ($(A)!{\k/6}!(C)$);
    \draw[thick] (M\k)--(N\k);
    \fill (M\k) circle (1.6pt) node[left]{\footnotesize$M_\k$};
    \fill (N\k) circle (1.6pt) node[right]{\footnotesize$N_\k$};}
    \foreach \p/\pos in {A/above,B/below left,C/below right,H/below}{\fill (\p) circle (2pt) node[\pos]{$\p$};}
\end{tikzpicture}
\end{center}
\begin{enumerate}
  \item Montrer que $M_kN_k=2k$ m pour chaque $k$.
  \item À quelle hauteur au-dessus de $[BC]$ se trouve la barre $[M_kN_k]$ ?
  \item Calculer la longueur totale des cinq barres.
  \item Pour $n$ barres qui partagent $[AB]$ en $n+1$ parties égales, on note $S$ la longueur totale. Montrer que $M_kN_k+M_{n+1-k}N_{n+1-k}=BC$, puis que $S=\dfrac n2\,BC$. [Écrire $S$ deux fois : dans l'ordre $k=1,\dots,n$, puis dans l'ordre inverse.]
  \item On dispose de $99$ m de barre pour un toit de base $BC=11$ m. Combien de barres peut-on poser au maximum ?
\end{enumerate}

% ---SOLUTION---
\begin{prop}[Thalès]
Si $M\in[AB]$, $N\in[AC]$ et $(MN)\parallel(BC)$, on part du sommet $A$ : \ib{$\dfrac{AM}{AB}=\dfrac{AN}{AC}=\dfrac{MN}{BC}$}.
\end{prop}
\begin{enumerate}
  \item Dans le triangle $ABC$, on a $M_k\in[AB]$, $N_k\in[AC]$ et $(M_kN_k)\parallel(BC)$. Thalès de sommet $A$ :
  \[\dfrac{AM_k}{AB}=\dfrac{AN_k}{AC}=\dfrac{M_kN_k}{BC}\]
  Or $\dfrac{AM_k}{AB}=\dfrac k6$ :
  \[M_kN_k=\dfrac k6\times12=\ib{$2k$}\ \text{m}\]
  Soit $2$, $4$, $6$, $8$ et $10$ m.

  \item La droite $(AH)$ est perpendiculaire à $(BC)$, donc à chaque barre. Elle coupe $[M_kN_k]$ en $K_k$.

  Dans le triangle $ABH$, on a $M_k\in[AB]$, $K_k\in[AH]$ et $(M_kK_k)\parallel(BH)$. Thalès de sommet $A$ :
  \[\dfrac{AM_k}{AB}=\dfrac{AK_k}{AH}\implies\dfrac{AK_k}6=\dfrac k6\implies AK_k=k\]
  La barre $[M_kN_k]$ est à la hauteur
  \[K_kH=AH-AK_k=\ib{$6-k$}\ \text{m}\]
  Les barres sont à $5$, $4$, $3$, $2$ et $1$ m du sol.

  \item On additionne les longueurs de la question 1 :
  \[2+4+6+8+10=\ib{$30$}\ \text{m}\]

  \item \textbf{Longueur d'une barre.} Même raisonnement qu'à la question 1, avec $AM_k=\dfrac k{n+1}\,AB$ :
  \[\ell_k=M_kN_k=\dfrac k{n+1}\,BC\]
  \textbf{Deux barres symétriques.} La barre numéro $n+1-k$ a pour longueur $\ell_{n+1-k}=\dfrac{n+1-k}{n+1}\,BC$. Donc :
  \[\ell_k+\ell_{n+1-k}=\dfrac{k+n+1-k}{n+1}\,BC=\ib{$BC$}\]
  \textbf{La somme.} On écrit $S$ dans les deux ordres :
  \[S=\ell_1+\ell_2+\cdots+\ell_n\]
  \[S=\ell_n+\ell_{n-1}+\cdots+\ell_1\]
  On additionne terme à terme : chaque paire $\ell_k+\ell_{n+1-k}$ vaut $BC$, et il y a $n$ paires.
  \[2S=n\times BC\implies\ib{$S=\dfrac n2\,BC$}\]
  Contrôle avec $n=5$ et $BC=12$ : $\dfrac52\times12=30$ m.

  \item La longueur totale pour $n$ barres est $S=\dfrac n2\times11$.
  \[\dfrac{11\,n}2\leq99\iff n\leq18\]
  On peut donc poser \ib{$18$ barres}, qui utilisent exactement $\dfrac{18}2\times11=99$ m.
\end{enumerate}
\textbf{Idée :} associer chaque barre à sa symétrique : leurs longueurs ont pour somme $BC$.

\textbf{Piège :} écrire $M_kN_k=\dfrac k6\,AB$ : le rapport de Thalès s'applique à $BC$, pas à $AB$.
